\subsection{Autofunzioni dell'oscillatore armonico}
\begin{align*}
a \ket{0} &= 0 \text{ è l'equazione che devo risolvere.} \\
\braket{x|n} &= \psi_n(x)
\end{align*}
\begin{align*}
\braket{x|a|0} = 0 &= \sqrt{\frac{m\omega}{2\hbar}} \braket{x|\left( \hat{x} + i\frac{\hat{p}}{m\omega} \right)|0} = \\
&= \sqrt{\frac{m\omega}{2\hbar}} \left( x \psi_0(x) + \frac{i}{m\omega} (-i) \hbar \frac{d}{dx}\psi_0(x) \right),
\end{align*}
\[
\frac{d}{dx} \psi_0(x) = -\frac{m\omega}{\hbar} x \psi_0(x),
\]
\[
\psi_0(x) = N_0 \exp \left\{ -\frac{m\omega}{\hbar} \frac{x^2}{2} \right\}.
\]
cioè è una gaussiana e rappresenta \textbf{lo stato fondamentale dell'atomo di idrogeno}.

Per gli stati superiori
\begin{align*}
\braket{x|1} = \psi_1(x) &= \braket{x|\sqrt{\frac{m\omega}{2\hbar}} \left( \hat{x} - i \frac{\hat{p}}{m\omega} \right)|0} = \\
&= N_1 \left( \exp \left\{ -\frac{m\omega}{\hbar} \frac{x^2}{2} \right\} - \frac{\hbar}{m\omega} \frac{d}{dx} \exp \left\{ -\frac{m\omega}{\hbar} \frac{x^2}{2} \right\} \right) = \\
&= N_1 \left( -\frac{\hbar}{m\omega} \right) \left( -\frac{m\omega}{\hbar} x \right) 2 \exp \left\{ \frac{m\omega}{\hbar} \frac{x^2}{2} \right\} = \\
&= \frac{N_1}{N_0} 2 x \psi_0(x).
\end{align*}
e poi iterando
\[
\psi_n(x) = \braket{x|n} = N_n \psi_0(x) H_n(x)
\]
dove \( H_n \) è un potenziale di grado \( n \). Sono i polinomi di Hermite.